\lhead{2. Boyer-Moore Algorithm}

\begin{enumerate}[label=\textbf{\Alph*.}]

\item \textbf{Define the Boyer-Moore algorithm. Why is it considered efficient for pattern
matching?}

The Boyer-Moore Algorithm is a powerful string matching algorithm used to search a pattern
inside a text. Unlike the Brute Force approach, it starts comparing characters from the
\textbf{rightmost side of the pattern}.

Whenever a mismatch occurs, it uses special rules to shift the pattern by multiple
positions instead of moving only one position at a time.

This makes the algorithm much faster and more efficient for large texts because many
unnecessary comparisons are avoided.

\item \textbf{Explain the following heuristics used in Boyer-Moore:\\
a) Bad Character Rule\\
b) Good Suffix Rule}

Heuristics used in Boyer-Moore:

\textbf{a) Bad Character Rule}

The Bad Character Rule is applied whenever a mismatch occurs.

If the mismatched text character appears somewhere earlier in the pattern,
the pattern is shifted so that both occurrences align.

If the mismatched character does not exist in the pattern,
the pattern is shifted completely beyond that character.

\textbf{Example:}\\
Text = ABCDXEF\\
Pattern = CDE

The character \textbf{X} causes a mismatch and is not present in the pattern.

Therefore, the pattern can be shifted several positions forward.

This reduces unnecessary comparisons.

\textbf{b) Good Suffix Rule}

The Good Suffix Rule is used when a portion of the pattern has already matched
successfully.

Instead of starting again, the algorithm looks for another occurrence of the
matched suffix within the pattern.

The pattern is then shifted accordingly.

This allows the algorithm to reuse previously matched information and improves\\
performance.

\item \textbf{Differentiate between Brute Force and Boyer-Moore algorithms.}

\begin{center}
\begin{tabular}{|c|>{\raggedright\arraybackslash}p{5.8cm}|>{\raggedright\arraybackslash}p{5.8cm}|}
\hline
& \centering\arraybackslash \textbf{Brute Force} & \centering\arraybackslash \textbf{Boyer-Moore}\\
\hline
1. & Compares from left to right & Compares from right to left\\
2. & Checks every position & Skips multiple position\\
3. & Shifts one position at a time & Uses larger shifts\\
4. & Easy to implement & Comparatively complex\\
5. & No optimization used & Uses optimization heuristics\\
\hline
\end{tabular}
\end{center}

\item \textbf{Apply Boyer-Moore Pattern Matching for the following:\\
\begin{tabular}{@{}rl@{}}
\quad Text: & THIS IS A TEST TEXT\\
\quad Pattern: & TEST
\end{tabular}
\\
Show the pattern shifts step by step.}

\textbf{Shift 1}

\begin{verbatim}
THIS IS A TEST TEXT
TEST
\end{verbatim}

Mismatch occurs.\\
Shift the pattern.

\textbf{Shift 2}

\begin{verbatim}
THIS IS A TEST TEXT
 TEST
\end{verbatim}

Mismatch occurs again.\\
Shift further.

\textbf{Shift 3}

\begin{verbatim}
THIS IS A TEST TEXT
      TEST
\end{verbatim}

Mismatch occurs.\\
Shift ahead.

\textbf{Final Shift}

\begin{verbatim}
THIS IS A TEST TEXT
          TEST
\end{verbatim}

Comparison from right to left:\\
T = T $\checkmark$\\
S = S $\checkmark$\\
E = E $\checkmark$\\
T = T $\checkmark$

All characters match.\\
Therefore, the pattern \textbf{TEST} is found at position \textbf{10}.

\end{enumerate}

\begin{center}
\vspace{1cm}
$\ast \ast \ast$
\end{center}
\newpage
